\subsection{1. Thiết kế kiến trúc hệ thống}
\subsubsection{1.1. Luồng xử lý tổng thể (System Pipeline)}
\indent\indent Kiến trúc hệ thống được thiết kế theo mô hình hai tầng liên kết chặt chẽ (Dual-stage Framework):
\begin{enumerate}[label=\textbf{Giai đoạn \arabic*:}, itemsep=3pt]
    \item \textbf{Phân loại đa nhãn dấu hiệu làm giả (Multi-label Classification):} Sử dụng mạng trích xuất đặc trưng \textbf{EfficientNet-B0} để phân tích ảnh vùng chữ nghi ngờ ($224 \times 224 \times 3$), dự đoán đồng thời 6 đặc trưng pháp y vi mô.
    \item \textbf{Trích xuất lập luận suy luận tự nhiên (Reasoning Extraction):} Tận dụng mô hình thị giác - ngôn ngữ \textbf{SmolVLM2-500M} để đọc vùng ảnh kết hợp cùng nhãn dự đoán từ Giai đoạn 1, sinh ra chuỗi tư duy giải thích chi tiết trong thẻ \texttt{<think>...</think>}.
\end{enumerate}

\subsubsection{1.2. Hệ thống nhãn phân loại pháp y (Forensic Labels)}
\indent\indent Dựa trên phân tích các hình thức làm giả văn bản trong thực tế và chuẩn gán nhãn của tập dữ liệu TFR/DTD, hệ thống định nghĩa 6 nhãn đặc trưng độc lập:
\begin{table}[H]
    \centering
    \caption{Đặc tả 6 nhãn phân loại dấu hiệu làm giả văn bản}
    \label{tab:labels_spec}
    \vspace{0.2cm}
    \renewcommand{\arraystretch}{1.3}
    \begin{tabular}{|m{0.8cm}|m{4.2cm}|m{9cm}|}
        \hline
        \textbf{STT} & \textbf{Tên nhãn (Label)} & \textbf{Mô tả đặc trưng vi mô} \\
        \hline
        1 & \textbf{Edge artifacts} & Vết cắt dán, viền răng cưa bất thường, dấu vết mờ nhạt do nội suy xung quanh nét chữ. \\
        \hline
        2 & \textbf{Unnatural texture} & Độ nhiễu nén JPEG, hạt nhiễu cảm biến hoặc độ nhẵn của giấy nền không đồng nhất với tài liệu gốc. \\
        \hline
        3 & \textbf{Inconsistent alignment} & Ký tự bị lệch dòng, khoảng cách giữa các chữ cái (kerning) hoặc khoảng cách dòng (leading) không chuẩn. \\
        \hline
        4 & \textbf{Inconsistent font} & Sử dụng font chữ khác biệt về độ dày nét (weight), độ nghiêng, hoặc kiểu dáng chân chữ (serif). \\
        \hline
        5 & \textbf{Lack of integration} & Vùng chữ thiếu tính hòa nhập vào nền, độ tương phản ánh sáng và đổ bóng không tự nhiên. \\
        \hline
        6 & \textbf{Text incoherence} & Nội dung văn bản bị vô nghĩa, mâu thuẫn ngữ nghĩa hoặc sai lệch logic với phần còn lại của văn bản. \\
        \hline
    \end{tabular}
\end{table}

\subsubsection{1.3. Cơ chế trích xuất chuỗi lập luận (Reasoning Rationale)}
\indent\indent Để cung cấp bằng chứng pháp y thuyết phục, SmolVLM2 được cấu hình theo cơ chế sinh văn bản tự hồi quy (autoregressive generation) với prompt có định dạng:
\begin{lstlisting}[language=Python, frame=single, caption={Định dạng Prompt suy luận cho SmolVLM2}]
messages = [
    {
        "role": "user",
        "content": [
            {"type": "image"},
            {"type": "text", "text": "Analyze the tampered region in this document. Explain the specific artifacts (font, edge, texture, alignment) and provide reasoning."}
        ]
    }
]
\end{lstlisting}
Mô hình sẽ sinh ra cấu trúc suy luận \texttt{<think> ... lập luận chi tiết ... </think>} trước khi đưa ra kết luận cuối cùng.
